% Hafta 3 — Kullanımda veri: açılma penceresini daraltan katmanlar
% Derleme: py -3.12 tools/latex/derle.py docs/week-3/assets/h03-18-kullanimda-veri-katmanlari.tex
\documentclass[tikz,border=8pt]{standalone}
\usepackage{cen429-cizim}
\begin{document}
\begin{tikzpicture}[node distance=6mm and 6mm]
  \node[baslik] (b) at (0,0) {Kullanımda veri: açılma penceresini daraltan katmanlar};
  \node[altbaslik] at (b.south west) {sır kullanılmak için bir an açılır; o an yönetilir};

  \node[iyikutu=48mm, minimum height=13mm, below=13mm of b.south west, anchor=north west, xshift=6mm] (k1)
    {\kb{Kısa ömür}};
  \node[font=\sffamily\small, text=soluk, text width=108mm, align=left, right=6mm of k1] (a1)
    {anahtarı işin hemen öncesinde aç, bitince sil (\kod{explicit\_bzero})};

  \node[kutu=48mm, minimum height=13mm, below=5mm of k1] (k2) {\kb{Bellek kilidi}};
  \node[font=\sffamily\small, text=soluk, text width=108mm, align=left, right=6mm of k2] (a2)
    {\kod{mlock} / \kod{VirtualLock} --- takas alanına yazılmasını engeller};

  \node[kutu=48mm, minimum height=13mm, below=5mm of k2] (k3) {\kb{Döküm kapatma}};
  \node[font=\sffamily\small, text=soluk, text width=108mm, align=left, right=6mm of k3] (a3)
    {core dump ve hata raporu devre dışı};

  \node[morkutu=48mm, minimum height=13mm, below=5mm of k3] (k4) {\kb{Cihaz bağlama}};
  \node[font=\sffamily\small, text=soluk, text width=108mm, align=left, right=6mm of k4] (a4)
    {anahtar cihazın kimliğine bağlanır; kopyalanırsa açılmaz};

  \node[iyikutu=48mm, minimum height=13mm, below=5mm of k4] (k5) {\kb{Donanım (TEE/SE)}};
  \node[font=\sffamily\small, text=soluk, text width=108mm, align=left, right=6mm of k5] (a5)
    {anahtar hiç uygulamaya girmez --- en güçlü katman};

  \node[dip, text width=175mm, below=10mm of k5.south, anchor=north]
    {Hiçbiri pencereyi sıfırlamaz; amaç saldırganın o ana yetişmesini pahalı kılmaktır.};
\end{tikzpicture}
\end{document}
